\section{Contoh Program}
\label{sec:contoh-bab8}

Dua program berikut dirakit dan dijalankan dengan Emu8086 atau
TASM/TLINK. Keduanya memperlihatkan bahwa mode pengalamatan langsung
menentukan bentuk kode dan alamat yang dibaca prosesor.

\subsection{Contoh 1: Membandingkan Tiga Mode pada Array}

Program \file{mode-array.asm} menjumlahkan lima elemen array yang
sama dengan tiga mode pengalamatan berbeda, lalu membandingkan ketiga
hasilnya. Kode lengkapnya tampak pada Kode~\ref{lst:mode-array}.

\lstinputlisting[
    language=MASM,
    caption={Program \file{mode-array.asm}: tiga mode pengalamatan
    untuk satu array.},
    label={lst:mode-array}
]{\CodeDir/mode-array.asm}

Cara kerja program dijelaskan sebagai berikut.

\begin{itemize}
    \item \textbf{Cara 1 --- mode direct.} Elemen dibaca dengan offset
          tetap, \lstasm|data1+0| sampai \lstasm|data1+8|. Mudah
          dibaca, tetapi tidak praktis untuk array panjang.
    \item \textbf{Cara 2 --- mode indexed.} \reg{SI} menjadi indeks
          yang bertambah dua setiap putaran; \lstasm|data1[SI]|
          setara dengan \lstasm|[data1+SI]| pada segmen \reg{DS}.
    \item \textbf{Cara 3 --- mode base-relative.} \reg{BX} diisi
          \lstasm|LEA BX, data1| sebagai basis, lalu elemen dibaca
          dengan \lstasm|[BX+SI]|; basis dan indeks terpisah sehingga
          mudah dipakai ulang.
    \item \textbf{Pembandingan hasil.} \asm{CMP} dan \asm{JNE}
          membandingkan \code{jumlah1} dengan \code{jumlah2}, lalu
          \code{jumlah2} dengan \code{jumlah3}. Struktur kontrol ini
          dibahas pada Bab~\ref{chap:kontrol} dan operasi
          aritmatikanya pada Bab~\ref{chap:aritmatika}.
\end{itemize}

Ketiga cara menghasilkan jumlah yang identik, yaitu
$10 + 20 + 30 + 40 + 50 = 150$ atau \hex{96}.

\begin{session}[caption={Keluaran program \file{mode-array.asm}.}, label={lst:keluaran-mode-array}]
Perbandingan tiga mode pengalamatan
Ketiga hasil identik.
\end{session}

\begin{infobox}
Bila hasil ketiga cara berbeda, kesalahan yang paling sering terjadi
adalah keliru memakai \reg{BP} sebagai basis. Bentuk \lstasm|[BP+SI]|
membaca segmen \reg{SS}, bukan \reg{DS}, sehingga data yang diambil
adalah isi stack, bukan isi array.
\end{infobox}

\subsection{Contoh 2: Membaca Parameter dengan \code{[BP+disp]}}

Program \file{param-bp.asm} memperlihatkan pemakaian mode
base-relative dengan register \reg{BP} untuk membaca dua nilai yang
diletakkan pada stack. Kode lengkapnya tampak pada
Kode~\ref{lst:param-bp}.

\lstinputlisting[
    language=MASM,
    caption={Program \file{param-bp.asm}: membaca parameter melalui
    \code{[BP+disp]}.},
    label={lst:param-bp}
]{\CodeDir/param-bp.asm}

Poin penting program ini adalah sebagai berikut.

\begin{enumerate}
    \item Nilai 2 dan 1 diletakkan pada stack melalui \asm{PUSH}.
          Prosesor 8086 tidak memiliki \lstasm|PUSH 2| untuk nilai
          segera, sehingga nilainya dimuat ke \reg{AX} terlebih
          dahulu.
    \item Setelah \lstasm|MOV BP, SP|, \reg{BP} menunjuk puncak
          stack. \lstasm|MOV AX, [BP+0]| mengambil parameter pertama
          dan \lstasm|MOV BX, [BP+2]| mengambil parameter kedua.
    \item Keduanya memakai segmen \reg{SS} secara bawaan. Bila
          penulisannya diubah menjadi \lstasm|[BX+2]| tanpa
          menyiapkan \reg{BX}, program akan membaca lokasi yang
          berbeda pada segmen \reg{DS}.
    \item Hasil penjumlahan bernilai 3 dan dicetak sebagai karakter
          melalui layanan DOS fungsi \hex{02}; konversi bilangan
          dibahas pada Bab~\ref{chap:aritmatika}.
    \item Dua instruksi \asm{POP} mengembalikan \reg{SP} ke posisi
          semula, kebiasaan yang akan dipakai secara baku pada
          Bab~\ref{chap:prosedur}.
\end{enumerate}

\begin{session}[caption={Keluaran program \file{param-bp.asm}.}, label={lst:keluaran-param-bp}]
1 + 2 = 3
\end{session}
